\documentclass[12pt,a4paper]{article}

% ---------- پکیج‌ها (همه قبل از xepersian) ----------
\usepackage{fontspec}
\usepackage{amsmath, amssymb, amstext}
\usepackage{graphicx}
\usepackage{geometry}
\usepackage{xcolor}
\usepackage{fancyhdr}

% ---------- تنظیم حاشیه ----------
\geometry{margin=1in}

% ---------- تنظیم fancyhdr ----------
\pagestyle{fancy}
\setlength{\headheight}{15pt} % حل خطای headheight

% هدرها
\lhead{کوییز درس مکانیک‌تحلیلی}
\chead{دانشکده فیزیک دانشگاه صنعتی خواجه نصیرالدین طوسی}
\rhead{\today}
\usepackage{xepersian}
\settextfont{B Nazanin}
\begin{document}
	
	% جدول اطلاعات فردی
	\begin{tabular}{|p{7cm}|p{7cm}|}
		\hline
		\textbf{نام و نام خانوادگی:} & \textbf{شماره دانشجویی:} \\
		\hline
	\end{tabular}
	
	% بخش سوالات
	\section*{سوالات}
	
	\begin{enumerate}
		\item
		یک خودرو با جرم \(1000\ \mathrm{kg}\) (به‌همراه سرنشینان) طوری است که به ازای اضافه شدن هر \(100\ \mathrm{kg}\) جرم، به اندازه \(1.0\ \mathrm{cm}\) بیشتر نسبت به سطح جاده نشست می‌کند.  
		خودرو با مولفه افقی سرعت ثابت \(20\ \mathrm{km/h}\) روی جاده‌ای حرکت می‌کند که به‌صورت یک موج سینوسی با برجستگی‌هایی به دامنه \(5.0\ \mathrm{cm}\) و طول موج \(20\ \mathrm{cm}\) مدل شده است. فاصله میان چرخ جلو و عقب \(2.4\ \mathrm{m}\) است.
		
		فرض کنید حرکت عمودی خودرو را می‌توان مانند یک نوسانگر هارمونیکِ تحتِ نیروی محرک و بدون میرایی مدل کرد. جرم چرخ‌ها و فنرها را نادیده بگیرید و فرض کنید چرخ‌ها همیشه با سطح جاده در تماس‌اند.  
		دامنه نوسان عمودی خودرو را به‌دست آورید.
		
		\item
		دو جرم با مقدار \(m_1 = 100\ \mathrm{g}\) و \(m_2 = 200\ \mathrm{g}\) روی یک مسیر افقی بدون اصطکاک به‌صورت آزاد حرکت می‌کنند و توسط یک فنر با ثابت نیروی \(k = 0.5\ \mathrm{N/m}\) به یکدیگر متصل شده‌اند.  
		بسامد نوسان این دستگاه را بیابید.
	\end{enumerate}
	
	\vspace{10pt}
	\textbf{موفق باشید.}
	
\end{document}
